% Vorlage für ein Inforz
% Erstellt von: Tobias Otterbein, totterbein@d120.de, Stand: 31.08.2014
% Darf beliebig für ein Inforz angepasst werden.
\documentclass[a5paper,pagesize,twoside,fontsize=10pt,DIV=15]{scrreprt}
\usepackage[utf8]{inputenc}
\usepackage{paket/inforz}
\usepackage{gensymb}
\usepackage{framed}
\usepackage[autostyle=true,german=quotes,english=american]{csquotes}
\usepackage{textcomp}
\usepackage{dashbox, anyfontsize}
%\usepackage{todonotes}
\usepackage{ marvosym }

% Sprache hier einstellen
\usepackage[ngerman]{babel}
\usepackage{blindtext}

% Auskommentieren, wenn keine Links im Inhalt benötigt werden. 4. Option sorgt für 2-seitige PDF-Ansicht
\usepackage[colorlinks=true,linkcolor=black,urlcolor=black,pdfpagelayout=TwoColumnRight]{hyperref}

% Schöne Anführungszeichen mit \enquote{[...]}
\usepackage{csquotes}

\begin{document}

    % Titelseite
    \begin{titlepage}~
\ThisCenterWallPaper{1}{grafik/cover}


\begin{textblock*}{12cm}(2cm,1.5cm)
\begin{flushleft}
\large\sffamily\textbf{
\newline
\textcolor{black}{Zeitschrift der Studierenden\\ 
der Informatik der \\
TU Darmstadt} } 
\end{flushleft}
\end{textblock*}


\begin{textblock*}{12cm}(1.5cm,1cm)
\includegraphics[width=12cm]{grafik/inforz}
\end{textblock*}


\begin{textblock*}{5cm}(1cm,7cm)
\begin{rotate}{90}
\sffamily\huge\textbf{
	\textcolor{black}{Inforz Sept. \the\year}}
\end{rotate}
\end{textblock*}


\begin{textblock*}{2cm}(1cm,14cm)
\begin{rotate}{90}
\textcolor{black}{\sffamily\small Preis: unbezahlbar}
\end{rotate}
\end{textblock*}


\begin{textblock*}{2cm}(1cm,20cm)
\begin{rotate}{90}
\textcolor{white}{\sffamily ISSN: 1614-4295}
\end{rotate}
\end{textblock*}

\end{titlepage}
\newpage


    % Vorwort
    \section*{Liebe Leser*innen,}

das Inforz ist ab jetzt komplett in Farbe gedruckt. Gerade rechtzeitig für unser tolles Cover und habt ihr schon rausgefunden, was das ist?
Solltet ihr das nicht wissen, schaut doch mal auf Seite ..., dort findet ihr einen Artikel, der sich damit beschäftigt.
Weiter haben wir tolle Artikel rund um Darmstadt und das Studium. So berichten wir über das HLMD (eines der komischen Gebäude neben der Uni) und erzählen euch ein wenig zum Arbeiten neben dem Studium.
\\\\
In der Freizeitrubrik stellen wir euch eine Auswahl an Podcasts vor und unserer Exchefredakteur präsentiert euch sein bestes Linsengericht.
\\\\
Auch freuen wir uns, dass wir zwei Interviews für euch geführt haben, zum einen mit Prof. Koch, den man aus EiCB, CMS und vielleicht sogar noch TGDI kennt. Das andere Interview haben wir mit Fr. Brandt geführt.
Auch wenn diese inzwischen in ihren Ruhestand gegangen ist, so war sie doch sehr viele Jahre die gute Seele der Fachstudienberatung. 
\\\\
Auch wieder dabei, sind die Filmkreistipps, sowie einige Linuxdistros und ein griechischer Buchstabe zum sammeln.
\\\\
Wir wünschen euch viel Spaß mit der neuen Ausgabe. Wenn ihr Anmerkungen habt oder selbst mal einen Artikel schreiben wollt, schickt uns einfach eine Mail an inforz@d120.de

\vspace{5mm}
Heiko Carrasco \& Jannis Blüml



\vfill


\bildmitunterschrift{grafik/comics/earth_orbital_diagram.png}{width=8cm}{}{xkcd.com/1878}





    % Inhaltsverzeichnis
    \tableofcontents
    \vfill
    \bildmitunterschrift{grafik/wesen/wesen_inforz}{width=2.5cm}{}{}
    \newpage

    % Inhalt

    %TODO HIER INHALT EINFÜGEN

    \phantomsection
    \addcontentsline{toc}{chapter}{Rund ums Studium}
    %Comic auf Seite 2 und/oder Seite 3
\artikel{Arbeiten neben dem Studium}{BaFöG, Stipendien, Eltern, Verwandte... Alle Studierenden benötigen für ihren Lebensunterhalt Geld und die meisten von uns haben eine Vielzahl von Quellen, aus denen sie ihre monatlichen Einkünfte beziehen. Immer wieder kommt es jedoch vor, dass das Geld, welches man bekommt, nicht ausreicht. Vielleicht reicht das Gehalt der Eltern nicht aus, um einen zu unterstützen, oder der Bafög Antrag wurde abgelehnt. Dann heißt es für viele: Arbeiten gehen.}
{   30 CP im Semester entsprechen in etwa einer Arbeitswoche von 40 Stunden. Da
    bleibt für viele gar nicht mal viel Zeit, auch noch einer geregelten
    Beschäftigung nachzugehen. Das erste, was ihr für euch entscheiden solltet,
    ist, wie viel Zeit ihr für eine Nebenbeschäftigung aufwenden könnt und
    wollt. Für die Uni, und ganz wichtig auch für den Staat, ist eure
    Hauptbeschäftigung das Studium, was rechtliche Konsequenzen hat. Mehr als
    20 Stunden die Woche während der Vorlesungszeit solltet ihr auf keinen Fall
    arbeiten, da ihr sonst sozialversicherungspflichtig werdet. Auch die
    Freibetragsgrenze (2016 waren es 8.472 Euro) solltet ihr im Auge behalten,
    da ihr sonst Steuern zahlen müsst.\\ \vfill \columnbreak Gut zusammengefasst bekommt ihr diese
    Informationen unter \url{https://www.studentenwerke.de/de/werkstudentenprivileg}. \\

    Prinzipiell sind eurer Kreativität keine Grenzen gesetzt, welchen Beschäftigungen ihr nachgehen wollt. Ob Nachtschicht an der Tankstelle, Werkstudent*in bei einem Softwareunternehmen oder selbstständiger Web-Designer, solange ihr die oben genannten Rahmenbedingungen einhaltet, könnt ihr arbeiten, wo ihr wollt\footnote{An Gesetze müsst ihr euch natürlich trotzdem halten, Drogen dealen ist nach wie vor verboten ;-)}. Allerdings unterscheiden sich unterschiedliche Formen der Arbeit voneinander! Die erste und wichtigste Unterscheidung ist, ob ihr angestellt oder selbstständig arbeiten wollt. \\

    Als Selbstständige seid ihr ziemlich frei, was die Gestaltung eurer
    Tätigkeit angeht, aber auf euch selbst gestellt. Als unabhängige
    Dienstleister könnt ihr eure eigene Firma gründen\footnote{Die  Gründung
        einer eigenen Firma ist ein Unterfangen, welches locker einen eigenen
        Artikel wert wäre. Um hier Platz zu sparen, verweise ich auf\\
        \url{https://www.existenzgruender-jungunternehmer.de/p/gruendung/studenten.html}
        für weitere Informationen.} und selbstständig eure Dienstleistung oder
    Produkte anbieten. Diese Freiheit heißt aber auch, dass ihr eine große
    Unsicherheit habt und viel Arbeit selbst machen müsst, denn Kunden
    anzuwerben und Preise zu verhandeln ist dann eure
    Aufgabe.\\\vfill\columnbreak Als Informatiker*innen gibt es aber sehr viele Nischen, in denen ihr euch sehr gut selbstständig machen könnt. Ob als Web Developer oder als Programmierer, die Grenzen setzen hier nur eure Zeit und euer unternehmerisches Talent. \\

    Als Angestellte seid ihr deutlich weniger frei, da ihr Arbeitnehmer*innen bei einem Arbeitgeber seid. Dafür könnt ihr dann auf in festes, vorhersehbares Gehalt freuen und müsst euch nicht darum kümmern, dass ihr Kunden bekommt\footnote{Es sei denn, ihr werdet zur Kundenwerbung eingestellt ;-)}. Außerdem seid ihr vor plötzlicher Kündigung geschützt und habt ein Anrecht auf Urlaub und Pausen. Euer größtes ''Problem'' dabei ist, einen passenden Arbeitsplatz zu finden. Dabei könnt ihr auf Ressourcen im Internet zurückgreifen, wie z.B. die Webseiten der Uni und das D120-Forum. Darüber hinaus lohnt es sich die Augen und Ohren offen zu halten. Oft suchen Freunde, Verwandte und Bekannte nach Mitarbeiter*innen oder wissen von interessanten Stellen. Und egal wie man zum berüchtigten ''Vitamin B'' steht, oft ist dies der beste Weg zu einem lukrativen und interessanten Job. \\

    Für Studierende steht ein Arbeitgeber natürlich immer offen: die Uni.
    Stellen als Hilfswissenschaftliche Mitarbeiter*innen (HiWis) sind
    Arbeitnehmerstellen. Eure Arbeitgeber*innen sind dann die verschiedenen
    Bereiche an der Uni, oft einzelne Fachgruppen oder die Fachbereiche.
    HiWi-Jobs sind zwar meist weniger gut bezahlt, als vergleichbare Stellen in
    der Wirtschaft (9,50\,\EURtm ~für Bachelor-, 11,50\,\EURtm ~für Master-Studierende), aber oftmals werden deutlich mehr Stunden bezahlt als man eigentlich leistet. Der größte Vorteil ist, dass die Arbeitgeber*innen an der Universität oft Verständnis für die Eigenheiten des Studiums haben. Hier erhaltet ihr oft mehr Nachsicht, wenn ihr wegen einer anstehenden Klausur weniger arbeitet, als in der freien Wirtschaft, wo die Produktivität des Unternehmens eindeutig Vorrang vor eurem Studium hat. \\

    Wichtig bei der Arbeitgeberwahl ist es, dass ihr euch darüber im Klaren seid, wer an eurer Arbeit mit verdient. An der Uni werben oft Vermittlungsfirmen, die euch an ihre eigenen Kunden weiterverleihen. Die Verträge mögen gut aussehen, aber oft fällt ein großer Teil eures eigentlichen Lohns an die Vermittlung. Sobald mehrere dieser Firmen hintereinander hängen, ihr also nur über mehrere Ecken weiter geleitet werdet, solltet ihr aufpassen. Es ist fast immer besser, direkt bei eurem entgültigen Arbeitgeber angestellt zu sein, als über Vermittlungsfirmen verliehen zu werden.Eine Außnahme bilden hier einige Consulting-Firmen, die euch ermöglichen, an interessante Projekte zu kommen, die ihr sonst niemals gesehen hättet. Die Entscheidung, ob die Vermittlungs- oder Consulting-Firma ihr Geld wert ist, müsst ihr je nach Fall selbst treffen. Wie an so vielen Stellen ist hier ein gesunder Menschenverstand gefragt. \\

    Die nächste wichtige Frage ist, wann ihr arbeiten wollt. Hier macht es vor allem einen Unterschied, ob ihr in den Semesterferien eher am Stück oder während der Vorlesungszeit nebenbei arbeiten wollt. Beides hat Vor- und Nachteile, aus persöhnlicher Erfahrung kann ich allerdings sagen, dass man in der Vorlesungsfreien Zeit zwar an sich viel Freiraum hat, aber ein großer Job das Lernen schnell in den Hintergrund rücken lässt. Seid also vorsichtig, wenn ihr euch für ein größeres Praktikum verpflichtet, bevor alle Klausuren geschrieben sind. Während der Vorlesungszeit habt ihr dafür natürlich auch andere Verpflichtungen, Vorlesungen, Übungen und andere Abgaben. \\

    Diese Übersicht soll nur einen schnellen Überblick über die wichtigsten
    Themen bei der Arbeitswahl bieten. Sollten euch bestimmte Themen genauer
    interessieren, schreibt uns eine kurze Mail an
    \texttt{\href{mailto:inforz@d120.de}{inforz@d120.de}}.

}
{Claas Voelcker}
\vfill
\bildmitunterschrift{grafik/comics/wizards.png}{width=11cm}{}{http://abstrusegoose.com/253}

    \clearpage
    \artikel{Ausflugstipps in Darmstadt: Das HLMD}
{Das Hessische Landesmuseum Darmstadt (HLMD) am Karolinenplatz wurde ab 2007 für sieben Jahre (zwei Jahre über Regelstudienzeit) und viel Geld (mehr als 300.000 Semesterbeiträge) komplett saniert. Nach der Wiedereröffnung im Herbst 2014 ist das Universalmuseum auf jeden Fall einen Besuch wert.
}{Und schon das erste Fremdwort. Ein kleiner Exkurs für Interessierte: Museum. Altgriechisch für Musentempel. Da nun nicht alle Expertenwissen in griechischer Geschichte haben, sei erwähnt, dass sie die Schutzgöttinen der Künste, Kultur und Wissenschaften waren. Ein Museum ist also ein Ort, an dem bedeutsame oder lehrreiche Gegenstände aus Kunst, Kultur und Wissenschaft aufbewahrt, erforscht und ausgestellt werden.

Also ein Museum, dass sich nicht auf einen Bereich wie Kunst oder Naturkunde spezialisiert, sondern wesentlich breiter aufgestellt ist. Dieses Konzept ist mittlerweile recht selten, was das HLMD in dieser Hinsicht besonders macht.

Aber was kann man dort denn nun alles sehen? Fangen wir von vorne an. Im Anfang schuf..., na gut so weit zurück dann doch nicht. Nur die letzten 50 Mio. Jahre. Fossilien aus allen möglichen Erdzeitaltern und Regionen können betrachtet werden. Zu den bekanntesten Fundstücken zählen sicherlich die Urpferdchen, die in der Grube Messel gefunden wurden. Und wenn man dann ein paar Millionen Jahre ins Land gehen lässt, landet man schon beim Mastodon, das freundlich aus dem 2. Stock ins Foyer blickt.

Mastodonten sind übrigens schon etwas länger ausgestorben. Und ja, auch ein Skelett kann freundlich gucken.

Kommen wir nun zur Zoologischen Sammlung, die viele interessante, faszinierende oder leider ausgestorbene Tiere und Vögel zeigt, sowie diverse Lebensräume als Dioramen mehr oder weniger realistisch darstellt. Es sei nochmal betont, dass sich die gezeigten riesigen Laufvögel weder als Haustier noch zum Reiten eignen. Außerdem sind sie -- richtig -- ausgestorben.

Genug Text über den Bereich Wissenschaft, der aber noch einiges mehr zu bieten hat. Nun folgt Kunst und Kultur. Auch hier gibt's eine breite Spanne von Epochen, Regionen und Künstler*innen zu sehen. Beginnend in der Steinzeit und im alten Ägypten geht es weiter nach Griechenland (schon wieder die Griechen) und in das römische Imperium. Uschebti, Vasen, ein Sarkophag, noch mehr Töpfe und Vasen, ein Mosaik, ein paar Säulen und einiges mehr. Zwischendrin dann noch ein paar Samurai. Die wollen schließlich auch gesehen werden.

Das finstere Mittelalter ist im HLMD gar nicht mehr so finster, sorgt doch der Catwalk mit den Ritterrüstungen für Glanz und Glamour, unterstützt durch handwerkliche Meisterwerke aus Glas, Holz und Metall.

Freunde der Malerei können sich auf Werke aus über 700 Jahren freuen, die im Anbau des Gebäudes ausgestellt werden. Und in diesem Fall dürften sich die meisten einig sein, dass das Kunst ist. Wer es moderner mag, muss den Anbau verlassen und unter das Dach des Hauptgebäudes steigen. Dort findet sich unter anderem die Müllhalde des Museums, oh Verzeihung, der Block Beuys. Kein weiterer Kommentar dazu. Moderne Kunst ist halt Geschmackssache. Ist doch völlig normal, dass man erstmal die Aufsicht fragen muss, ob hier noch aufgebaut wird oder nicht.

Ach ja, Jugendstil gibts auch in hinreichend großer Menge zu sehen. Wir sind ja schließlich in Darmstadt, einer der Hauptstädte des Jugendstils.

Zusammenfassend lässt sich sagen, dass das HLMD definitiv einen Besuch wert ist. Mit 4 Euro ist der Eintritt für Studierende recht günstig und Eintrittskarten eigenen sich auch recht gut als Weihnachtsgeschenk für Eltern o.ä. Die Erfahrung hat gezeigt, dass Viele hunderte Kilometer in andere Städte und Länder reisen, um Dinge zu sehen und zu erleben, aber irgendwie nie auf die Idee kommen, die eigene Stadt und die nähere Umgebung zu erkunden. Das ist recht schade, denn auch die (temporäre) Heimat bietet viele Dinge, die man nicht erwarten würde. Das HLMD liegt quasi immer auf dem Weg zur Uni, doch drin waren sicherlich die wenigsten. Daher ein Appell: Erkundet Darmstadt und seine Umgebung, denn auch hier kann man tolle / faszinierende / interessante Dinge entdecken.}{Tobias Otterbein}

\bildmitunterschrift{grafik/HLMD.jpg}{width=\textwidth}{Hessisches Landesmuseum Darmstadt}{Jan Bambach}

    \clearpage
    \definecolor{myBlue}{RGB}{188,219,238} %www.farb-tabelle.de/en/table-of-color : LightSteelBlue2


\newcommand\addHeader[3]{
    \noindent
    \begin{tabular}{lp{3cm}}
        \noindent Titel & #1 \\
        \scriptsize Originaltitel & \scriptsize #2 \\ 
       % \scriptsize Filmkreistermin & \scriptsize #3
    \end{tabular}\\
    \vspace{0.125em}\\
}
\newcommand\addBeschreibung[1]{#1\\}
\newcommand\addAnmerkung[1]{}%\vspace{0.4em}\\\begin{tabular}{lp{3cm}} \small Anmerkungen &
%\small \textit{#1}\end{tabular}}
\newcommand\addOV{}%\addAnmerkung{Originalvertonung}}
\newcommand\addOMU{}%\addAnmerkung{Originalvertonung mit Untertiteln}}
\newcommand{\film}[1]{\parindent-\fboxsep\par\fbox{\parbox{\linewidth}{#1}}\\}

\artikel{Mein Filmtipp zum Filmkreisprogramm}
{Unter dem Titel "`The greatest program ever! It's huge!"' brachte der Filmkreis auch dieses Semester wieder viele (über 50 Filme) tolle und ausgezeichnete Filme auf die Wand des Audimax. Ich habe mir das Programm (größtenteils) angeschaut und mir einige Filme herausgesucht, die ich empfehlen würde und über die ich kurz reden möchte.}
{
    \film{
        \addHeader{Fantastische Tierwesen und wo sie zu finden sind}{Fantastic Beasts and where to find them}{02.03.2017}
    Was soll ich sagen, ich bin mit Harry Potter aufgewachsen und liebe es noch heute. Es gab keinen Zweifel, dass dieser Film es in meine Highlights schaffen wird. Er dient als Auftakt zu einer ganz neuen Reihe und verspricht viel. Beeindruckend ist auch, dass ich das Gefühl hatte, dass der Film mit seinem Publikum geht und sich mehr an die Kinder von damals (hust hust) richtet, als an die neue Generation. 
    \addOV
}
\film{
    \addHeader{Zwei Missionare}{Porgi l'altra guancia}{16.03.2017}
    Als Kind durfte ich immer mit meinem Vater die alten Terence Hill und Bud Spencer Filme sehen und irgendwie haben sie mich verzaubert, auch wenn ich sie damals hasste. Zwei Missionare ist ein gelungener Film der beiden "`Kultgestalten"' und lässt mich in Erinnerungen schwelgen. Wer Filme der beiden kennt, denen muss ich nicht erzählen mit was sie rechnen müssen.An schönes Andenken an den verstorbenen Bud Spencer
} 
\film{
    \addHeader{Die Melodie des Meeres}{Song of the Sea}{22.03.2017}
    Animationsfilm? Kenne ich! Verschwundene Mutter? Kenne ich! Märchenhaft? Kenne ich. Als ich das erste Mal auf eine Beschreibung dieses Films traf, klang es für mich eher wie einer dieser klassischen Kinderfilme und schrieb ihn als solchen ab. Was für ein Fehler von mir. Klar ist der Film für Kinder und auch die Handlung nicht die neuste, aber die Bilder und Animationen machen das alles wett. Wer in Zeiten von Disney und Pixar einen Blick in einen anderen Artstyle werfen möchte, der freut sich, dass dieser Film es in den Filmkreis geschafft hatte.
    \addOMU
}
\film{
    \addHeader{Kubo - Der tapfere Samurai}{Kubo and the Two Strings}{13.04.2017}
    Da wir schon bei außergewöhnlichen Animationsfilmen sind, zeigte der Filmkreis noch einen weiteren sehr Empfehlenswerten. Kubo ist das neueste Werk des Studios "`Laika"', die sich auf Stop Motion Filme spezialisiert haben. Ich empfehle jedem mal ein "`Behind the scenes"' von Laika zu schauen, erst dann begreift man den Aufwand, der bei solchen Filmen anfällt. Dabei ist Kubo ein wunderschöner Film geworden, der seine Zuschauer auf seine ganz eigene Weise verzaubert. 
    \addOV
}
\film{
    \addHeader{Rogue One}{Rogue One: A Star Wars Story}{20.04.2017}
    Star Wars ist ähnlich wie Harry Potter etwas, mit dem ich aufgewachsen bin. Rogue One führt uns wieder in diese unglaubliche Welt und zeigt uns abseits von Lichtschwertkämpfen und der Geschichte rund um den Auserwählten einen kleinen Teil seiner Vielseitigkeit. Dabei darf natürlich ein Todesstern nicht fehlen.
    \addOV
}
\film{
    \addHeader{Die glorreichen Sieben}{The Magnificent Seven}{27.04.2017}
    Die Neuauflage eines Klassikers, der so alt ist, dass ich ihn nicht mehr kenne. Ich werde also keinen Bezug zum Original nehmen. Der Western ist gut gemacht, lädt mit guten Schauspielern und reichlich Action zum Zuschauen ein. Die Handlung ist denkbar simpel, was den Film aber in keinster Weise stört, liegt sein Ziel sowieso eher auf Action und dem Aufbau eines großen Showdowns.
}
\film{
    \addHeader{1984}{Nineteen Eighty-Four}{16.05.2017}
    Wer kennt ihn nicht, den Klassiker von George Orwell. Für mich klar eine Pflichtlektüre für jeden Informatiker, diese Umsetzung des Films aus dem Jahr 1984 zeigt die wesentlichen Aspekte und erfasst das Bedrückende und schwere Thema des Buches gut. Ich kann jedem empfehlen diesen Film mal zu sehen, gerade weil ich nicht das Gefühl habe, als hätten die Themen an Bedeutung verloren, ganz im Gegenteil.
    \addOV
}
\film{
    \addHeader{Snowden}{Snowden}{18.05.2017}
    Wem dieser Name nichts sagt, der lebt wohl hinterm Mond. Trotzdem heißt kennen noch lange nicht, wissen was alles passiert ist. Dieses Biopic zeigt die Geschehnisse rund um die NSA Affäre und gerade Joseph Gordon Levitt in seiner Rolle sticht für mich heraus. Sollte man diesen Film nicht sehen wollen, so empfehle ich Citizenfour. Einen der Beiden sollte man aber gesehen haben. 
}
\film{
    \addHeader{Spiel mit das Lied vom Tod}{}{08.06.2017}
    Euer Wunschfilm des Semesters! Nicht ohne Grund wurde dieser zu eurem Lieblingsfilm gekrönt. Ein durchausgelungener Film, wenn auch nicht umbedingt mein Genre.
    \addAnmerkung{Was gezeigt wird liegt an euch!!!}
}
\film{
    \addHeader{Captain Fantastic - Einmal Wildnis und zurück}{Captain Fantastic}{06.07.2017}
    Ich glaube, für mich das Highlight dieses Filmprogramms! Viggo Mortensen in einer, wenn nicht seiner besten Rolle und ja, ich weiß, dass er Aragorn gespielt hat. Der Film zeigt uns ein so weltfremdes Szenario, das trotzdem begeistert und mich fasziniert zurücklässt.
}
\film{
    \addHeader{Arrival}{Arrival}{18.07.2017}
    Aliens kommen auf die Erde und die Menschen starten einen Krieg. Dieses Szenario kennen wir alle zur Genüge, aber was, wenn die Menschen keinen Krieg anfangen? Dieser Film beschäftigt sich mit genau einem solchen Szenario, dass großartig erzählt ist.
}
\film{
    \addHeader{Der Schuh des Manitu - Extra large}{Der Schuh des Manitu - Extra large}{20.07.2017}
    Also bitte... \\
    \\
    \textit{You don't have to wait for later 
    here's a new eliminator \\
    Ask your local weapon-trader\\ 
    for the Superperforator...\\
\\}
    ...Muss ich mehr sagen? 

}
}
{Jannis Blüml}

	\vfill
	\bildmitunterschrift{grafik/comics/homosapiensapien.png}{width=\textwidth}{}{http://abstrusegoose.com/572}


    \clearpage
    \artikel{Sprüche und Zitate}
{"'Als ich an diese Uni kam, da gab es in den Inforzen noch coole Sprüche und Zitate von Professoren, wo sind die denn hin?"' (ein Leser)}
{
    Ganz einfach, wir können solche Zitate nur einbringen, wenn sie uns mehr mitgeteilt werden. Wir schaffen es nicht in jeder Vorlesung zu sitzen, weshalb uns viele dieser Sprüche durch die Lappen gehen. Wenn euch allerdings diese Rubrik gefallen hat, so schreibt doch einfach an inforz@d120.de und schickt uns eure besten Sprüche, die ihr von Professor*innen aufgeschnappt habt. Nur so können wir wieder welche abdrucken. \\

    \bildmitunterschrift{grafik/wesen_wa(h)l.png}{width=5cm}{}{}
}
{die Inforzredaktion}
\dictum[Prof. Koch, EiCB]{"`Strings sind böse, das sollen Sie als wesentliche
    Erkenntnis dieser Lehrveranstaltung mitnehmen."' (Prof. Koch, EiCB)}
\vspace{1em}~\\
Als während einer Frage das Telefon im Publikum ertönte:\vspace{.5em}
\dictum[Prof. Mantel, FMIS]{"`No, don't call your friends in the lecture. This is not one of these
    shows, where you can ask a friend for an answer."'}

\bildmitunterschrift{grafik/comics/fun_with_statistics.png}{width=11cm}{}{http://abstrusegoose.com/54}

    \clearpage
    \phantomsection
    \addcontentsline{toc}{chapter}{Aus der Fachschaft}
    %Comic am Ende oder Anfang von anderem Artikel auf gleicher Seite
\artikelohnevorspann{Podcasts für Informatiker*innen}
{Die meisten Menschen stecken, wenn sie nach der Arbeit die Uni verlassen, sich erstmal Kophörer rein. Ich bin da keine Ausnahme. Aber während die meisten Anderen dann direkt in den
    Musikabspieler ihrer Wahl wechseln, schalte ich meine Podcastapp ein und verbringe die nächsten Minuten oder auch Stunden damit, dem Geplapper anderer Leute zuzuhören.
    Aber was ist denn ein Podcast? Im Prinzip ist es meist eine oder mehrere Personen, die sich über ein Thema unterhalten. Beispielsweise gibt es Podcasts über Filme und Serien oder welche, die sich mit Technologie und Wissenschaft beschäftigen. Dabei gibt es innerhalb der Casts mehrere Episoden, ähnlich einer Serie. Diese sind nicht unbedingt zusammenhängend, man kann also auch in der Mitte einsteigen. Die Länge variiert, liegt aber oft im Bereich von einer Stunde und mehr.
    Mit der Zeit haben sich im Internet eine Fülle an Podcasts angesammelt, sodass die Auswahl am Anfang recht schwer fällt. Daher möchte ich hier eine kleine Auswahl meiner abonnierten Podcasts vorstellen. Die Liste ist vor allem wissenschaftlich/technisch orientiert. Alle vorgestellten Casts sind kostenlos, über Spenden (z.B. mit flattr)
    freuen sich die meisten. Sie können über einen beliebigen Podcatcher (Podcastapp) bezogen werden. Unter Android ist zum Beispiel PodcastAddict sehr zu empfehlen.
    \vfill\columnbreak
    ~\vfill
    \textbf{Name: }Binärgewitter\\
    \textbf{Sprache: }Deutsch\\
    \textbf{Thema: }OpenSource- und Techniknews\\
    \textbf{Beschreibung: }In Binärgewitter treffen sich ein paar alte Hasen der Podcastszene und quatschen über das aktuelle Geschehen in der Welt der Technik.
    Aktuelle News werden vorgestellt und kommentiert. Dabei nehmen sich die Macher selbst nicht unbedingt ernst und oft enden Diskussionen in dem ein oder anderen
    Schlagabtausch.
    \vfill\vfill
    \textbf{Name: }Chaosradio\\
    \textbf{Sprache: }Deutsch\\
    \textbf{Thema: }Computer und Gesellschaft, Technik, CCC\\
    \textbf{Beschreibung: }Das Chaosradio kommt jeden Monat raus. Vom CCC Berlin organisiert, werden vor allem Themen welche die Hackerszene betreffen besprochen. Das kann zum Beispiel Netzpolitik, Biohacking oder digitales Geld sein. Dazu werden oft Experten eingeladen, zum Beispiel Leute von Netzpolitik.org oder der FSFE\footnote{\textit{Free Software Foundation Europe}}. Abgerundet wird die Sache von der guten Moderation, durch einen professionellen Radiomoderatoren.
    \vfill~\\\columnbreak\\
    \textbf{Name: }Serial\\
    \textbf{Sprache: }Englisch\\
    \textbf{Thema: }Die Geschichte eines Mordes\\
    \textbf{Beschreibung: }In Serial erzählt Sarah Koenig die wahre Geschichte über den Mord an Hae Min Lee im Jahre 1999 in Baltimore, USA. Im Laufe der Geschichte deckt die Reporterin Ungereimtheiten und
    Fehler während der Aufklärung und Verhandlung des Falls auf. Jede Episode befasst sich mit einem anderen Teilaspekt des Falls, wodurch sich nach und nach ein Gesamtbild zusammenpuzzelt.
    \vfill\columnbreak
    \textbf{Name: }Robot Congress\\
    \textbf{Sprache: }Englisch\\
    \textbf{Thema: }Technologie und Recht\\
    \textbf{Beschreibung: } Dieser Podcast hat sich der juristischen Einschätzung von technischen Themen verschrieben. Mit Ryan Morrison, besser bekannt als der Video Game Attorney, ist auch jemand vom Fach an Bord. Trockene juristischen Texte muss man aber nicht fürchten, die Themen werden mit Humor angegangen. Generell bekommt man aber leider nur die amerikanische Seite des Rechts zu sehen.}{Heiko Carrasco}

    \vspace{1em}
    \artikel{Rezepte aus der Fachschaft: \\
	Linsencurry à la Totterbein}
{Informatiker*innen und Kochen, das kann nicht klappen. So das weitverbreitete Vorurteil. Also was machen?}
{Hier wieder ein Rezept zum nach kochen. Diesmal ein veganes Linsencurry. \\

	\textbf{Zutaten \\ (im praktischem 4er-WG Maßstab)}
	\begin{itemize}
		\item 2 Paprika (rot)
		\item 2 Zwiebeln (mittel-groß)
		\item 1 Dose Kichererbsen
		\item 2 Chilis
		\item 4 Knoblauchzehen
		\item 2 Karotten
		\item 200\,g Linsen (rote)
		\item 500\,ml Kokosmilch
		\item ggf. weitere 200\,ml Kokosmilch
		\item 2\,EL Currypaste (rote)
		\item ca. 1\,EL Tandoori Masala
		\item ca. 1\,EL Kurkuma
		\item eine Prise Salz und Pfeffer
		\item 250\,g Basmatireis
	\end{itemize}
	\vfill~\columnbreak\\
	\textbf{Zubereitung} \\
	Man sollte anfangen die Linsen in heißem, nicht kochendem Wasser einzulegen und aufzuweichen.
	Diese brauchen bei mir ungefähr eine Stunde, bis sie weiterverarbeitet werden können. In dieser Zeit kann man anfangen die Zwiebeln, Paprika, Chilis, Karotte und den Knoblauch zu schneiden. Ich schneide hier immer recht kleine Würfel, an sich ist die Form und Größe aber Geschmackssache. Die Zwiebeln sind mit dem Knoblauch in einer Pfanne bei mittlerer Hitze zu dünsten. Sobald diese glasig sind, können wir die Paprika, Karotte und Chilis dazugeben. Alles nochmal kurz anbraten. In der Zwischenzeit kann man Tomaten vierteln und anschließend dazugeben. Nun mit Salz, Pfeffer, Kurkuma und Tandoori würzen und abschmecken. Der nächste Schritt ist es, die Zwiebeln, den Sellerie und die Karotten zu schälen (sofern nötig) und in Würfel zu schneiden. Sobald diese fertig geschnitten sind, kann man in einer Pfanne die Butter langsam mit etwas Öl(nur ein wenig) erhitzen, dies sollte allerdings nur auf geringer Temperatur stattfinden. Sobald die Butter geschmolzen ist, gibt man Zwiebeln, Sellerie und Karotten in die Pfanne und lässt diese 30-45 Minuten auf dem Herd. Dabei zwischendurch immer mal wieder umrühren. Man kann nun die Chilis schneiden und entkernen, den Knoblauch pressen oder kleinschneiden und die Tomaten würfeln, dann kann man in einer zweiten Pfanne das Hackfleisch scharf anbraten. Sobald dieses durch ist, löscht man es mit einem Schuss Rotwein ab und gibt die Milch dazu. Ist der Karotten-Sellerie-Zwiebeln-Mix fertig, so kann man diesen mit dem Hackfleisch zusammen in einen großen Topf geben, den restlichen Rotwein dazugeben. Dies wird die spätere Soße. Hier gibt man jetzt noch Tomaten, Chili und Knoblauch dazu und lässt die Soße 3-5 Stunden auf kleiner Temperatur köcheln. Die Soße wird so reduziert, sollte sie zu stark reduzieren, dann gebe die Gemüsebrühe hinzu. Dabei sollte immer mal wieder umgerührt werden. Zum Ende der Zeit fängt man an, die Nudeln zu kochen und fertig ist das Essen. Am Ende nicht das Würzen vergessen. Alles in eine feuerfeste Form oder Topf geben und mit den Linsen und Kichererbsen vermengen. Danach Kokosmilch und Currypaste dazugeben, Deckel drauf. Alles bei 160°C (Umluft) 2 Stunden im Ofen lassen. Immer mal wieder reinschauen und bei zu viel Flüssigkeitsverlust etwas Kokosmilch dazugeben. Anschließend nochmal nachwürzen und fertig ist das Curry. Zu diesem empfehle ich als Beilage Basmatireis oder Naan.
}
{Jannis Blueml}
\clearpage
\vfill

%	\bildmitunterschrift{grafik/rezept_flowchart.png}{width=1.0\textwidth}{}{Kevin Otto}

\vfill
\newpage


    \clearpage
    \artikelgruppe{Linuxdistributionen Teil 2}
{Nach dem wir letztes Mal bereits angefangen haben einige unser Lieblingsdistributionen vorzustellen, wollen wir euch hier nochmal zwei Alternativen präsentieren. Elementary OS, das Mac unter den Linuxdistros, und Mint, das Windows unter ihnen. Wer eine objektive Meinung erwartet, der erwartet zu viel. Hier eine rein subjektive Meinung einiger Individuen. }
{
	\subartikelohnevorspann{elementaryOS}
	{Elementary OS ist ein auf Ubuntu basierendes Betriebssystem, dass sich vor allem an Personen wendet, die mehr auf "`Eyecandy"' setzen, als auf Funktionalität. Dabei orientiert sich das System sehr stark an OSX/macOS. Das Betriebssystem ist leichtgewichtig und kommt mit seinem eigenen Desktop Environment (Pantheon), welches die Ziele der Entwickler perfekt umsetzt. Die Grundidee der Entwickler hinter ihrem Design ist nämlich wie folgt definiert: "`Sei prägnant, vermeide Konfiguration und gebe minimale Dokumentation"', sprich es soll intuitiv sein, soll "`Out of the Box"' funktionieren und gut vorbereitet sein und es soll verständlich sein, sodass man als Benutzer keine große Dokumentation braucht. \\ 
		
	Das versuchen die Entwickler auch möglichst gut umzusetzten, so setzen sie
    auf eine einfache Bedienoberfläche und einige Standardprogramme. \linebreak Leider sind nicht alle davon fehlerfrei oder zu 100 Prozent ausgearbeitet und durchdacht. So gibt es zum einen viele Bugs, die mit der Zeit die Nutzung stören oder erschweren, aber auch einige Designentscheidungen die ich nicht nachvollziehen kann. Hierbei fallen mir unter anderem die vordefinierte Dockingbar ohne File Browser ein, so wie das Blockieren des Systems bei der Installationen von Software, die nicht in ihrem "`AppCenter"' angeboten werden (ist etwas umständlich möglich). Auch wenn die Entwickler ein "`Out of the box"'-Feeling geben wollen, so sollten sie mich als Nutzer nicht einschränken, was die Konfiguration und Anpassung des Systems an mich betrifft. Dies passiert allerdings viel zu häufig und stört mich dann doch wieder. \\

	Von seiner Schwierigkeit ist es sehr leicht zu bedienen und eignet sich somit auch für Anfänger, wenn ich dort auch Mint vorziehen würde, gerade wenn ich mir die Fehleranfälligkeit von Elementary anschaue. \\
	
	Was das Betriebssystem möchte ist sicherlich sinnvoll, viele Nutzer wünschen sich ja schließlich ein "`schönes Linux"', welches wenig aufwand macht, leicht zu bedienen ist und nicht viel Hintergrundwissen erfordert. All das verspricht einem Elementary OS und mit etwas mehr Zeit wird es das eventuell auch schaffen, doch in seinem jetzigen Zustand ist es noch keine Alternative. \\	
	Denn so schön die Oberfläche auch ist, so vermisse ich viele Funktionen, die mir andere Distributionen wie Mint, Arch oder Fedora zur Verfügung stellen. Auch ist es mir noch zu instabil und hat zu viele Fehler. Ich sehe aber Potenzial in der Idee. }
	{Jannis Blüml}
	\subartikelohnevorspann{Linux Mint}
	{Es gehört zu den beliebtesten Linux der Zeit und ist auf Platz eins der Distrowatch Liste. Durch Projekte wie den Cinnamon Desktop und
	die Nutzung der Ubuntu Paketquellen hat sich Mint schnell als benutzerfreundlicher Ersatz für Canonicals Betriebssystem erwiesen
	und viele Menschen zum Umsteigen gebracht. In diesem kurzen Test möchte ich zeigen, warum das so ist und für wen sich diese Distro lohnt.\\
	Mint ist für Linuxverhältnisse recht alt. Bereits 2006 wurde die erste Version veröffentlicht. Schon damals wurde auf die Ubuntu Paketquellen
	gesetzt, was relativ aktuelle Software und etablierte Tools wie die Paketverwaltung apt und dem bekannten Installer 
	dem/der Nutzer*in zur Verfügung stellt. Auch gibt sich das Projekt viel Mühe, eine vollständige Kompatibilität mit Ubuntu aufrecht zu erhalten.
	Dadurch sind viele Anleitungen und Programme auf beide Distros
	äquivalent anwendbar.\\
	Einen großen Unterschied gibt es jedoch: Die angebotenen Desktopumgebungen. Hier steht bei Mint nämlich nicht Unity sondern standardmäßig
	Cinnamon, eine Eigenentwicklung, und Mate zur Verfügung. Zudem werden XFCE und KDE angeboten, aber meist mit größerer Verzögerung.
	Cinnamon ist ein aufgeräumter Desktop, welcher sich an Windows orientiert. Als Fork von GNOME wurde versucht, ein klassisches Bedienkonzept
	zu fahren, welches dem Nutzer möglichst nicht im Weg steht, aber trotzdem einiges an Konfigurationsmöglichkeiten bietet. Beim Testen erwies
	sich der Desktop als sehr schnell und elegant, hat aber standardmäßig die nervige Angewohnheit, bestimmte Aktionen des Nutzers mit Tönen zu
	untermalen. Die lässt sich jedoch in den aufgeräumten Einstellungen ausschalten. Auch viele weitere Möglichkeiten der Konfiguration gibt es
	hier, Cinnamon lässt sich ganz den Wünschen seiner Nutzer*innen anpassen.\\
	Die mitgelieferte Software entspricht den Standardprogrammen der meisten Distros. Firefox als Browser und Thunderbird als Mailprogramm
	kommen standardmäßig zum Einsatz. Zudem wird LibreOffice mitgeliefert,
    sowie eine Auswahl an Hilfsprogrammen wie Transmission und Gimp.\\
	Insgesamt handelt es sich größtenteils um FOSS\footnote{Free Open Source Software}, jedoch bietet der Installer an,
	unfreie Plugins und Codecs wie z.B. Flash nachzuladen und zu
    installieren.\columnbreak\\
	Bezüglich Updates hat sich das Mintteam leider für eine recht harsche Politik entschieden. Sowohl das Betriebssystem als auch die
	installierten Programme werden nicht geupdatet, wenn keine hunderprozentige Kompatibilität gewährleistet ist. Dadurch wurden in der
	Vergangenheit bereits Sicherheitsupdates erst mit großer Verspätung ausgeliefert. Andererseits macht das Mint gerade für Anwender*innen
	attraktiv, welche hohe Stabilität brauchen. Der daraus teilweise
    resultierende Sicherheitsverlust lässt sich dann eventuell verschmerzen.
	\medskip\\
	Linux Mint ist eine sehr schöne und stabile Distribution. Der robuste Kern und der gut konfigurierbare Desktop machen diese Linuxversion
	attraktiv für eine große Breite an Anwender*innen, sowohl Anfänger*innen als auch Expert*innen. Das Betriebssystem wirkt fertig und durchdacht,
	die Grundfunktionen sind schnell erreichbar und laufen sehr flüssig, auch auf älteren Rechnern. Die teilweise etwas veraltete Software
	lässt sich verschmerzen, die Updatepolitik ist hingegen eher fragwürdig. Insgesamt ist Mint jedoch absolut empfehlenswert und definitiv für mich
	eine Installation wert.
	}
	{Heiko Carrasco}}
	
	\vfill
	\bildmitunterschrift{grafik/comics/screenshots.png}{width=\textwidth/4*3}{}{xkcd.com/1863}


\newpage

    \artikel{Esoterische Programmiersprachen - Teil 1: Piet}
{Dies ist der erste Artikel einer Reihe über esoterische Programmiersprachen.\\
Aber was ist das eigentlich?\\
Esoterische
Programmiersprachen sind nicht zum praktischen Gebrauch bestimmt, die Ziele
mit denen diese entwickelt werden sind aber recht unterschiedlich.\\
Bei der
Programmiersprache \glqq{}Brainfuck\grqq{}, die vermutlich recht viele kennen, war die Prämisse
eine turingvollständige Sprache zu entwickeln, die mit einem möglichst kleinen
Compiler auskommt. \glqq{}Whitespace\grqq{} degegen soll beim Ausdruck des
Sourcecodes möglichst wenig Tinte bzw. Toner verbrauchen.}
{Die Programmiersprache, die ich heute vorstelle heißt \glqq{}Piet\grqq{}, benannt nach dem
Künstler Piet Mondrian. Ihr Ziel ist es gängigen Mondrian Bildern (s.
Abb 1) möglichst nahe zu kommen.
\bildmitunterschrift{grafik/piet_bild.jpg}{width=.75\columnwidth}{\small{}Composition
with Red, Yellow and Blue}{Piet Mondrian}
 Ihr fragt euch vielleicht, wie eine Programmiersprache (bzw. ein in ihr geschriebenes Programm) einem abstrakten
Gemälde ähneln soll.
Nun, in Piet geschriebene Programme sind ebenfalls Bilder.
Abbildung 2 zeigt ein Piet Programm, welches das Wort \glqq{}Piet\grqq{} ausgibt und auf
dem Umschlag ist ein Programm abgebildet, welches das Wort \glqq{}Inforz\grqq{}
erzeugt.
\bildmitunterschrift{grafik/piet-program.png}{width=.75\columnwidth}{\glqq{}Piet\grqq}{Thomas
Schoch}~\\

Im nachfolgenden werde ich die grundlegenden Konzepte erläutern:\\

Ein Programm wird aus den Farben Rot, Gelb, Grün, Cyan, Blau und Magenta
  in jeweils drei Helligkeitsstufen Hell, Normal und Dunkel, sowie aus den
  Farben Weiß und Schwarz zusammengesetzt.\\

Ein zusammenhängender Block von Pixeln einer Farbe wird \glqq{}Colour Block\grqq{}
  genannt.\\

Zur Programm Ausführung werden zwei Zeiger verwendet: Der \glqq{}Direction Pointer\grqq{}
  (DP)
  zeigt nach oben, unten, links oder rechts, während der \glqq{}Codel Chooser\grqq{} (CC) nach
  links oder nach rechts zeigt.
  Die Ausführung beginnt im Block in der oberen linken Ecke mit DP=rechts und
  CC=links
  Für einen Ausführungsschritt wird nun zunächst die Kante des aktuellen
  Blocks gewählt, welche sich am weitesten in Richtung des DP befindet.
  Nun wird das Pixel ausgewählt, dass sich an dieser Kante am weitesten in
  richtung CC (\grqq{}aus sicht des DP\grqq{}) befindet von dort wählt der Interpreter den
  anliegenden Block, führt die Instruktion dieses Überganges aus und führt den
  nächsten Schritt aus.\\

Instruktionen sind die Übergänge zwischen zwei verschiedenfarbigen Blöcken.
  So wird zum Beispiel bei einem Übergang zu einer um eine Stufe dunkleren
  Farbe (im Kontext von Piet ist \glqq{}Hell\grqq{} eine Stufe dunkler als \glqq{}Dunkel\grqq{}, so dass
  ein Zyklus entsteht) die Größe des Blockes von dem der Interpreter kommt auf den
  Stack gepusht.
  Bei einem Übergang zur nächsten Farbe (Reihenfolge wie oben, nach Magenta
  kommt wieder Rot) werden die letzten beiden Werte vom Stack genommen addiert
  und die Summe wieder auf den Stack gepusht.\vspace{5em}\\

Wenn der Interpreter bei der Ausführung einen schwarzen Block betreten
  müsste, oder wenn das Bild an dieser Seite zu Ende ist wird statt den neuen
  Block auszuwählen zunächst der CC geändert und vom aktuellen Block aus erneut
  versucht einen Schritt auszuführen. Falls dies auch schief geht, wird der DP
  im Urzeigersinn gedreht. Diese beiden Schritte werden so oft wiederholt, bis
  die Programmusführung fortgesetzt werden kann, oder bis der Interpreter
  wieder im Ursprungszustand ist (nach insgesammt 8 Versuchen). In diesem Fall
  wird das Programm beendet.\\

Wenn der Interpreter einen weißen Block betritt, so wird dieser einfach
  passiert und die Ausführung in gerader Linie am nächsten farbigen Block
  fortgesetzt. Weder das Betreten noch das Verlassen des weißen Blockes führt
  irgendeine Instruktion aus.\\

~\\~\\
Insgesammt ist Piet eine recht amüsante Abwechslung. Allerdings braucht man
schon für die Entwicklung kleiner Programme unverhältnissmäßig viel Zeit
weshalb die Programmiersprache für den alltäglichen Gebrauch eher ungeeignet
ist.\\
Wer sich dennoch mal an der Programmierung in Piet versuchen möchte dem
empfehle ich sich auf \url{http://www.dangermouse.net/esoteric/piet.html}
umzuschauen.

}
{Fabian Franke}

    \clearpage
    \artikel{Archiv: Griechische Buchstaben}
{Wir gehen in die sechste Folge fächerübergreifend beliebten Sammelserie der griechischen Buchstaben.
Schnell ausscheiden und ins Sammelalbum einkleben oder mit Freunden tauschen!}
{\vspace*{-1.5\baselineskip}\begin{center} 
	\includegraphics[scale=0.4]{grafik/zeta.png}
\end{center}
\flushleft
Wir verlassen die seichten Gebilde der Anfängerbuchstaben und 
nähern uns den etwas unbekannteren Buchstaben, die nicht so
häufig verwendet werden. Aber trotzdem sollte in der vollständigen
Sammlung der Griechischen Buchstaben das $\zeta$ nicht fehlen,
schon alleine, um etwas humanistische Bildung zu demonstrieren.

\subsection*{Verwendung}

Es gab mal einen berühmten Mathematiker namens Riemann, der hat eine Funktion
erfunden und diese $\zeta$-Funktion genannt. Sie stellt die Verteilung
der Primzahlen als Formel dar. Und es gab auch mal ein paar Informatiker, die
haben ein Betriebssystem mit diesem Namen versehen, das früher auch mals BeOS
bekannt war.
Der Autobauer Lancia hat ein Auto namens $\zeta$ herausgebracht.
Viel mehr Beispiele der Verwendung gibt es eigentlich nicht, da dieser
Buchstabe so unbekannt ist, verwenden ihn nur wenige.

\subsection*{Zubereitung}
Man nehme einen Stift, zum Beispiel den Lieblingsfüller
mit gefederter und ergonomisch geformter Grundmulde.
Mit diesem nun das Papier berühren und dabei die Bewegungen
ausführen, die zum Erscheinen des $\zeta$ auf selbigen
führen. Diese wären: eine kleine Schleife gegen den Uhrzeigersinn,
dann eine große hintenangehängt. Wenn eine halbe
Umdrehung vollbracht ist, die Richtung ändern und einen kleinen Schlenker
unten anhängen. Fertig ist das $\zeta$ und kann fortan bewundert werden.

\subsection*{Empfehlung}
Da der geübte Laie bei der Zubereitung des $\zeta$ leicht eine
Sehnenscheidenentzündung bekommen kann, empfehlen wir, eine ausgewogene
Mischung aus $\zeta$ und anderen Buchstaben zu verwenden. Sieht dann auch nicht so
eintönig aus.}
{Arne Pottharst}\vfill
\small Dieser Artikel erschien ursprünglich im Juni 2006 und steht daher nicht unter
CC-BY-SA.

    \clearpage
    \phantomsection
    \addcontentsline{toc}{chapter}{Interviews mit unserem Fachbereich}
    % !TeX spellcheck = de_DE


\artikel{Interview mit Professor Koch}
{Als Professor Koch vor zwölf Jahren an die Uni kam und das Fachgebiet ESA gründete, interviewten wir ihn zum ersten Mal. Dieses Interview könnt ihr in der Ausgabe April 2005 finden. Nun wollten wir wissen, wie sich der Fachbereich und das Fachgebiet ESA entwickelt haben.
    \bildmitunterschrift{grafik/koch}{width=\linewidth}{Prof. Koch im Interview April 2017}{}
}
{
    \subsubsection*{Warum sind Sie Professor geworden?}
    Aus zwei Gründen, einmal bin ich an der Wissenschaft und speziell an der Ingenieurstätigkeit sehr interessiert, andererseits lehre ich auch sehr gerne. Ich habe damit recht früh angefangen und habe schon in meiner Schulzeit meinen Mitschülern und Lehrern C beigebracht. Aus diesem Grund erschien mir die Tätigkeit als Professor passender als zum Beispiel Wissenschaftler in einem Forschungsinstitut zu werden.

    \subsubsection{War die Informatik bzw. die Elektrotechnik schon immer Ihr Ziel?}
    Ich habe mich schon sehr früh für Technik und Elektronik interessiert und habe damals Kurse für Elektronik an unserer Volkshochschule belegt. Irgendwann kam dann aber der Punkt, an dem mir mathematische und physikalische Grundlagen gefehlt haben. Auch hatte ich großen Respekt vor diesen Transistoren, denn wenn man damit rumspielte tendierten sie dazu, kaputt, zu gehen und das war nicht gut für mein Taschengeld. Deswegen war die Elektronik für mich nur so was wie eine Einstiegsdroge. Als Anfang der 80er die Computer aufkamen, war ich sofort begeistert, denn erstens konnte man durch Programmieren nichts kaputt machen und zweitens war man nicht durch den Zwang des Aufbaus neuer Hardware behindert. Dieser neuartige Prozess, von der Idee zur fertigen Umsetzung, ohne beispielsweise einen Lötkolben in die Hand nehmen zu müssen, hat mich so gefesselt, dass ich mich schlussendlich für die Informatik entschied.

    \subsubsection{Sie sind der einzige Professor dieses Fachbereichs, der sich so speziell mit Hardware beschäftigt. Fehlen uns hier Kompetenzen und sollten wir in Zukunft in diesen Bereich mehr investieren?}
    Das ist eine der wenigen Dinge, die ich in den letzten Jahren bedauere. Der Fachbereich hat einige sehr attraktive Forschungsgebiete unterstützt, dabei ist leider die technische Informatik stark geschrumpft. Als ich hier angefangen habe, gab es noch drei Professuren in diesem Bereich. Zwar tun wir unser Möglichstes, ein gewisses Lehrangebot aufrecht zu erhalten, aber es gibt auch einige Vorlesungen, die ich aus Zeitgründen schon seit Jahren nicht mehr halten konnte. Ich könnte mir vorstellen, gerade in Zusammenhang mit sehr aktuellen Themen, diesen Bereich wieder aktiver zu gestalten, zum Beispiel im Bereich "`\textit{High Performance Computing}"'. Hier erreichen wir so langsam die Grenzen der bisherigen Ansätze. Es gibt aber frische Forschungsideen, die sich mit neuartigen (Rechner-)Architekturen beschäftigen, um einige aktuelle Probleme anders zu lösen. Hier wäre eine gute Zusammenarbeit zwischen den Anwendern und den "`Hardwareaffinen"' notwendig, um eine möglichst gute Optimierung zu finden. Allerdings sehe ich auch das Problem der begrenzten Professorenstellen, man kann nicht alle Bereiche gleich gut ausbauen. Dass es in den letzten Jahren die technische Informatik aber so stark getroffen hat, finde ich persönlich sehr schade.

    \subsubsection{Sie haben vor knapp zwölf Jahren das Fachgebiet ESA gegründet. Hat es sich seitdem so entwickelt, wie Sie es sich erhofft haben?}
    Inzwischen ja, obwohl wir immense Schwierigkeiten in der Aufbauphase hatten. Ich habe damals ja nur mit einem Mitarbeiter angefangen. Und da wir doch recht ingenieursnah arbeiten, war die Arbeit recht aufwendig und wir litten doch unter unserer unterkritischen Masse. Mit der Zeit erkannte der Fachbereich allerdings die überhöhte Arbeitsbelastung auch in der Lehre und hat für Unterstützung gesorgt. Mit der aktuellen Größe von 8-10 Leuten bin ich recht zufrieden, auch wenn wir zwischenzeitlich mal etwas größer waren. Da hatte ich aber dann das Gefühl, dass ich meine Doktoranden nicht mehr intensiv genug betreuen konnte und ich zu weit weg von ihrer Arbeit war. Ich bin nicht Professor geworden, um nur Manager zu sein, sondern ich möchte mich mit meinem Fachgebiet auch unmittelbar befassen. Mit der aktuellen Größe bin ich sehr zufrieden, da ich dadurch nah genug an meinen Mitarbeitern bin, um auch fachlich mitzuhalten und mitzulernen.

    \subsubsection{Im aktuellen Aufbau der Vorlesung \textit{Architektur und Entwurf von Rechnersystemen} (AER) lernt man "`Bluespec Systemverilog (BSV)"' oder kurz Bluespec kennen. Was ist denn aktuell Ihre Lieblings Hardware Description Language?}
    Das ist eine sehr gute Frage, allerdings muss man hier mit "`Für welchen Zweck?"' noch präzisieren. Wenn es darum geht, hoch produktiv und höchst optimiert zu arbeiten, dann ist Bluespec die beste HDL, die wir zurzeit zur Verfügung haben. Aber auch Bluespec ist nicht ganz ohne: Durch seine Abstraktion ist es nicht immer ganz leicht zu verwenden und man benötigt einiges an zusätzlichem Wissen. Deswegen wäre es am besten, wenn man Hardware für spezifische Anwendungsgebiete in der Sprache beschreiben könnte, die in dem jeweiligen Anwendungsgebiet Verwendung findet. Hätten wir beispielsweise eine Perl-artige Sprache, die speziell in der Netzsicherheit verwendet wird, so wäre ein Compiler von z.B. dem relevanten Perl-Subset zu BSV sinnvoll, wenn sich die Sicherheitsexperten nicht mit Bluespec oder Verilog auskennen. Hier entwickelt man sogenannte Domain Specific Languages (DSL): Diese orientieren sich dann an den natürlichen Beschreibungen ihrer Anwendungsdomäne und  schaffen es, daraus Hardware zu erzeugen. Dies ist natürlich der eleganteste Weg, Hardware zu erstellen. Solche DSLs werden daher in Zukunft noch eine große Rolle spielen.

    \subsubsection{Sie halten aktuell regelmäßig die Veranstaltungen AER, EiCB, RO und DT. Diese vier Veranstaltungen liegen allerdings alle im Pflichtbereich. Was wären denn Veranstaltungen, die Sie gerne mal wieder halten würden, die Sie zeitlich aber aktuell nicht schaffen?}
    Ich würde gerne eine Veranstaltung zum Thema Moderne Rechnerarchitekturen anbieten, die in unserem Curriculum derzeit fehlt. In RO schauen wir uns MIPS an, in AER bekommen die Studierenden einen kleinen Einblick in ARM. In den letzten Jahren ist das Gebiet aber stark gewachsen. Wenn ich Zeit hätte würde ich gerne eine Vorlesung vorbereiten, in der wir uns genau diese neuen Trends einmal anschauen. Dabei sollten vom kleinen embedded Prozessor über Multi-Core CPUs und Grafikkarten alle Architekturen und Geräte wenigstens übersichtsartig behandelt werden. Es ist wichtig, die Funktion dieser Recheneinheiten auch auf Hardware-Ebene zu kennen, um sie besser programmieren zu können. Leider fehlt mir hierfür gerade einfach die Zeit.

    \subsubsection{In der Hoffnung, wir behalten Sie weitere zwölf Jahre, was wären Ihre Pläne für diese?}
    Das hängt ganz davon ab, wie sehr ich in Zukunft in Pflichtaufgaben eingebunden werde. Wenn ich weiter, wie in den letzten Jahren, so viele Pflichtveranstaltungen betreue, dann bremst das doch schon sehr aus. Sowohl in Bezug auf weiterführende Lehrveranstaltungen, als auch in der Forschung. Man kann leider nicht alles gleichzeitig machen. Ein konkretes Ziel ist, in Zukunft enger zusammen mit den Anwendungsexperten aus speziellen Domänen an DSLs zu arbeiten und zu forschen. Aus diesen DSLs kann man dann anwendungsspezifische Hardware generieren, die für ihre Domäne so optimiert ist, dass sie herkömmliche Prozessoren in den Schatten stellen kann. Eine Domäne, in der das aktuell schon sehr gut funktioniert, ist das Maschinelle Lernen. Ich denke aber, da geht noch mehr!

    \subsubsection{Wenn Sie sich die letzten zwölf Jahre anschauen, was waren denn Ihre Highlights in Ihrer Zeit als Studiendekan?}
    Zweifellos die Reakkreditierung, speziell das Auseinandersetzen mit unserem Curriculum und seine Weiterentwicklung, damit es heutigen Anforderungen genügt. Dabei musste ich aber auch im Auge behalten, dass es Studierenden und Dozenten gleichermaßen liegt.

    \subsubsection{Was ist denn Ihr Highlight aus den kompletten 12 Jahren?}
    Singuläre Ereignisse wie die Reakkreditierung gibt es eigentlich wenige, höchstens noch den Erhalt des Athenepreises für gute Lehre. Das hat mich sehr gefreut und ich werte das ganz hoch. Wir freuen uns aber natürlich auch, wenn wir unsere Forschungsergebnisse in guten Tagungen oder Journals unterbringen.

    \subsubsection{Wenn wir uns mal kurz etwas von der Uni entfernen, was sind denn Ihre Hobbys außerhalb Ihres Berufs?}
    Leider erschreckend wenig. Ich bin ein spielender Mensch, das kann man schon sagen. Dabei interessiert mich weniger das Glücksspiel, wohl aber Brett- und Computerspiele. Gerade letztere haben mich auch zur Informatik selber gebracht. Ich war doch fasziniert, wie so ein Atari VCS funktionierte. Ich probiere gerne vieles aus und schaue, wie es mir gefällt. Ich gehe gerne Wandern und lese viel, wenn ich Zeit habe. Hier dann idealerweise keine Fachbücher, sondern eher fantastische Literatur.

    \subsubsection{Was ist denn Ihre Lieblingsklausurfrage?}
    "Was sind die Stufen des Pipeline MIPS Prozessors, die wir kennen gelernt haben?". Die stelle ich zum Aufwärmen immer wieder gerne.

    \subsubsection{Und was ist Ihre Lieblingsantwort?}
    Ich habe am liebsten Antworten, bei denen der Studierende nachgedacht und nicht nur aus der Vorlesung reproduziert hat. Das kommt leider viel zu selten vor.
    Gerade wenn ein Studierender die Antwort nicht kennt, sollte er versuchen, die Antwort selbst zu entwickeln und die richtigen Schlüsse zu ziehen, anstatt die Frage einfach auszulassen.

    \subsubsection{Im letzten Interview sagten Sie, Ihre Lieblingsprogrammiersprache sei Smalltalk. Sehen Sie dass immer noch so?}
    Das ist immer noch richtig, aber die Sprache ist leider tot. Nicht, weil sie schlecht ist, sondern weil die Anbieter von Entwicklungswerkzeugen und ähnlichem mit ihrer Lizenzpolitik die Sprache geschlachtet haben. Die Eleganz, die Smalltalk hatte, habe ich bisher noch in keiner anderen Sprache gefunden. Aktuell oben auf meiner Liste von noch zu lernenden Sprachen stehen aber Scala und Haskell.

    \subsubsection{Was verwenden Sie, Windows, Linux oder Mac?}
    Alle drei. Es kommt darauf immer an, was man machen will. Ich bin aber persönlich ein Freund von Unix. Wenn ich die Wahl habe, verwende ich Linux, aber wir wissen alle, dass das nicht immer ganz so gut klappt.

    \subsubsection{Was ist Ihr liebster Texteditor / TeX-Editor?}
    Als Editor für Texte verwende ich Kile,  wobei ich aktuell auf Texstudio wechsle. Das liegt daran, dass Kile Texte so umformatiert, dass die Funktionalität von Git o.ä. etwas behindert wird.

    \subsubsection{Was ist für Sie die 42?}
    Die Antwort auf die all bewegende Frage, die Deep Thought gefunden hat und für mich, neben der 23, die Standardkonstante für alle meine Vorlesungen.

    \subsubsection{Warum die 23?}
    Ich weiß es ehrlich gesagt nicht, das habe ich mir von anderen Kollegen angewöhnt. Allerdings kenne ich die Geschichte dahinter nicht.\footnote{Die Zahl geht wohl auf die {\em Illuminatus}-Romanreihe von Shea und Wilson zurück.}

    \subsubsection{Letztes mal sagten Sie, Sie haben keine Zeit zum Kochen. Hat sich das inzwischen geändert?}
    Ich bemühe mich und werde auch immer besser.

    \subsubsection{Fehlt Ihnen irgendwas an den aktuellen Inforzen?}
    Es gab früher immer eine Seite mit Glossen meiner Kollegen und von mir, die fand ich immer sehr unterhaltsam.

    \subsubsection{Was wollen Sie den Studenten noch mitteilen?}
    Sie sollen keine Scheu haben zu denken. Rufen sie nicht nur die gecachten Antworten aus dem Gedächtnis ab, sondern verstehen sie diese auch.

}
{Heiko Carrasco \& Jannis Blüml}

\vfill
\begin{minipage}{0.5\textwidth}
    \bildmitunterschrift{grafik/comics/quantum.png}{width=\textwidth}{}{xkcd.com/1861}
\end{minipage}
\begin{minipage}{0.5\textwidth}
    \bildmitunterschrift{grafik/comics/physics_confession.png}{width=\textwidth}{}{xkcd.com/1867}
\end{minipage}

\newpage

    \clearpage
    \artikel{Interview mit Ulrike Brandt}
{Ulrike Brandt war lange eine Konstante am Fachbereich Informatik. Viele Studierende kennen sie und haben bei ihr Antworten auf viele Fragen über die Formalia des Studiums gefunden. Wir wollten von Frau Brandt wissen, was sie aus den Jahren an der Uni mitnimmt und was ihre Pläne für den Ruhestand sind.
}
{
\subsection*{Sie sind ja inzwischen eine ganze Weile hier, aber war denn die Informatik immer Ihr Ziel?}
Nein, eigentlich wollte ich in die Naturwissenschaften und so was wie Mathematik oder Chemie studieren. 
Als ich mir hier in Darmstadt einen Studiengang in Richtung chemische Verfahrenstechnik angeschaut habe, bin ich quasi über die Informatik gestolpert. Da habe ich mir gedacht, ich probiere das mal aus. Es klang halt ganz interessant.
\subsection*{Wie war denn das Informatikstudium zu ihrer Zeit?}
Bei mir war es so, dass ich nach meinem Vordiplom nicht wusste was die Informatik eigentlich ist. Damals gab es noch keinen roten Faden. Es gab zwar sehr interessante Vorlesungen, von zum Beispiel Robert Piloty, Hans-Jürgen Hoffman und Hermann Walter, aber die Zusammenhänge sind mir nicht immer ganz klar geworden. Wäre nicht die theoretische Informatik gewesen, die mir von Anfang an sofort Spaß gemacht hat, hätte ich hier wahrscheinlich die Informatik an den Nagel gehängt.
\subsection*{Sie sind nach dem Studium beziehungsweise nach Ihrer Habilitation in die Industrie gegangen. Was haben Sie da getan?}
Ich muss sagen, ich war nur ungefähr drei Jahre weg. In denen habe ich für die Höchst AG mathematisch-technische Assistenten ausgebildet. Danach habe ich kurz in der Firma meines Vaters als Geschäftsführerin gearbeitet, bis ich glücklicherweise wieder hierher zurückkommen konnte.
\subsection*{Wollten sie denn schon immer an der Uni arbeiten oder ergab sich das einfach so?}
Ja, da der Bereich, welcher mich am meisten interessierte, schon immer die theoretische Informatik war und dies auch der Bereich ist, in dem ich damals mehr machen wollte. Daher war die Universität eigentlich schon immer mein Ziel. Hier ist man dafür am besten aufgehoben. Es war Glück, dass das alles so gut geklappt hat und ich hierher zurückkehren konnte, ansonsten wäre ich halt in der Industrie geblieben.
\subsection*{Sie haben ja damals noch Diplom studiert und kennen deswegen sowohl dieses als auch das neuerer Bachelor/Master System. Was ist denn ihre Meinung zum neuen System bzw. zu dem Wechsel weg vom Diplom?}
Aus meiner ganz persönlichen Sicht, wenn ich das entscheiden könnte, wäre ich wahrscheinlich beim Diplom geblieben. Ich bin der Meinung, dass das Bachelor/Master System stärker reguliert ist und als solches einige Freiheiten verloren gehen, die es damals gab. Auch wenn der Master bei uns schon sehr viele Wahlmöglichkeiten hat und damit in die richtige Richtung geht. Ebenso gab es im Diplom weniger Prüfungen, zum Beispiel vor dem Vordiplom nur 5 Prüfung und einige "`Scheine"', die man erbringen musste. Vom Umfang her würde ich sagen, dass eine solche Prüfung von damals den Aufwand von zwei Fächern heute widerspiegelt, aber trotzdem ist das dann immer noch sehr viel weniger als dies heute der Fall ist. Gleiches gilt natürlich auch für das Hauptdiplom.
\subsubsection*{Was würden Sie sagen, waren ihre schönsten Erlebnisse in knapp 25 Jahren an der TU?}
Der Abschied war sehr nett, den fand ich sehr herzlich.
\subsubsection*{Und zur Gegenfrage, was würden Sie denn am liebsten wieder vergessen?}
Also so schlecht, dass ich es vergessen möchte ist es nicht, aber was ich sehr bedauert habe, ist die Abschaffung der Theorie. Sie ist zwar nicht komplett abgeschafft und in vielen Veranstaltungen integriert, allerdings wurde sie schon sehr stark zusammengekürzt. Dadurch das Aussagen und Prädikatenlogik (APL) und Automaten, formale Sprachen und Entscheidbarkeit (AFSE) sehr früh im Studium kommen, sind diese Veranstaltungen nicht besonders tief. Viele Bereiche, die ich als sehr wichtig einschätze, bleiben leider etwas auf der Strecke. Das bedauere ich doch schon sehr. 
\subsubsection*{Wir haben Sie bereits vor 9neun Jahren im Rahmen der Fachstudienberatung interviewt (Ausgabe: Oktober 08). Was hat sich an Ihrer Arbeit in dieser Zeit verändert?}
Ich mache inzwischen weniger Studienberatung, das hat Tim [Neubacher, Anm. der Redaktion] zu großen Teilen übernommen. Ich habe mich seitdem verstärkt um die Weiterentwicklung des Fachbereichs und der Studiengänge gekümmert. So war ich beispielsweise ziemlich stark in der Reakkreditierung vor zwei Jahren involviert.
\subsubsection*{Damals (vor neun Jahren) sagte man uns: "`Jeder Student sollte seine Prüfungs- und Studienordnung mal gelesen haben!"'. Hat sich daran etwas geändert oder vertreten sie immer noch diese Meinung?}
Ja, auf jeden Fall! Man sollte sie zumindest mal gelesen haben und die rudimentären Sachen mitgenommen haben. Man sollte auch immer wissen, an wen man sich wenden kann, sollte man mal einen Rat brauchen.
\subsubsection*{Was würden sie neuen Studierenden an dieser Uni raten?}
Wichtig ist, in einer Gruppe zu arbeiten und sich nicht zu isolieren. Auch sollte man sich immer über sein Studium informieren und sich notfalls helfen lassen. Einfach "`vor sich hin zu studieren"', ohne sich damit auseinandergesetzt zu haben ist einfach nicht gut. Nehmt Angebote wie die Ophase war und geht lieber einmal zu oft, als einmal zu wenig, zur Studienberatung.
\subsubsection*{Sie saßen ziemlich lange in Gremien wie dem LuSt. Wie fanden sie ihre Arbeit dort, konnten Sie so mitwirken wie sie es sich gewünscht hätten?}
Teils, teils. Das ist halt so. Im Ganzen bin ich aber recht zufrieden. Während der Reakkreditierung 2003 war der Kommission der Anteil der Mathematik in unserem Studiengang zu viel, weshalb wir die Mathematik  mit der Elektrotechnik Mathematik zusammengelegt haben. Das fand ich damals nicht so gut, hier hätte ich mir gewünscht, dass es anders gelaufen wäre. Aber alles in allem konnte ich vieles einbringen und bin sehr zufrieden damit.
\subsubsection*{Sie haben sehr viel Zeit nicht nur in unserem Fachbereich, sondern insgesamt an dieser Uni verbracht. Wie hat sich denn die Uni im Laufe ihrer Zeit hier entwickelt?}
Also ich glaube, dass die Professoren mehr unter Stress stehen, weil das Einwerben von Drittmitteln hier inzwischen eine große Rolle eingenommen hat. Was sich dadurch zum Beispiel geändert hat, ist, dass früher die Professoren, und zwar alle, im Fachbereichsrat saßen und sich mehr in die Lehre als Ganzes eingebracht haben. Damals waren wir auch nicht so viele Professoren. Ich war bereits in der Fachschaft, damals nannten wir uns noch Basisgruppe, und saß als Vertreterin lange und auch oft im Fachbereichsrat. Und das Engagement jedes einzelnen Professors sich einzubringen war viel mehr vorhanden. Heute ist es größtenteils das (Studien-) Dekanat und einige wenige Professoren, die sich als Gruppe versuchen einzubringen. Ich wünsche mir, das sich dies irgendwann wieder ändert und mehr zurück zum Zustand wie es früher einmal war zurückgeht.
\subsubsection*{Wenn Sie irgendwas an diesem Fachbereich ändern könnten, egal was, was wäre das?}
Die stärkere Beteiligung und mehr Interesse an Lehre, denn nur so können wir den Fachbereich besser gestalten.
\subsubsection*{Haben sie eine konkrete Idee, wie man die Professoren mehr einbeziehen könnte?}
Ich glaube, auch wenn das jetzt etwas pessimistisch klingt, es ist aktuell hoffnungslos. Ich sehe auch die Problematik nicht am mangelnden Interesse mehr an der Lehre mitzuwirken, sondern mehr darin, dass sie keine Zeit haben sich mehr einzubringen.
\subsubsection*{Könnten sie sich vorstellen, dass eine Art "`Vollversammlung aller Professoren, Mitarbeitern und der Fachschaft"' helfen würde? }
Man könnte das mal versuchen, aber ob das zum gewünschten Ergebnis führt weiß ich nicht. Vielleicht wäre es aber eine Idee, so etwas mal zu machen und dann zu evaluieren.
\subsubsection*{Sie haben ja demnächst etwas mehr Zeit als vorher. Wissen sie schon, was sie mit Ihrer neuen Freizeit anfangen werden?}
Also teils, teils. Erstmal mache ich ein bisschen Urlaub. Ich habe von meinen Kollegen zum Abschied einen Gutschein für einen Garten geschenkt bekommen und werde diesem wohl demnächst etwas Zeit zu Gute kommen lassen. Ich mach erstmal das, was mir Spaß macht. Dann sind in letzter Zeit einige Dinge liegen geblieben, denen ich mich wieder etwas mehr widmen möchte. Zum Beispiel kam meine wissenschaftliche Arbeit etwas zu kurz. Ich habe ein fast fertiges Paper, dass ich demnächst angehen möchte. Vielleicht will ich auch noch eine Ausbildung in Richtung Psychologie machen, dass hat mich schon immer interessiert. Ich hatte sogar überlegt es zu studieren.
\subsubsection*{Was nehmen sie aus 25 Jahren Uni mit und wie hat es sie verändert?}
Man braucht eine ganze Menge Ausdauer und Zähigkeit, wenn man hier arbeitet. Manchmal laufen die Dinge nicht immer einfach oder so wie man es sich vorstellt. Das braucht Kraft. Ich habe hier gelernt, dass man an bestimmten Stellen halt nicht einfach mit dem Kopf durch die Wand kann.
\subsubsection*{Windows, Mac oder Linux}
Ich habe mich an Windows gewöhnt und nie etwas anderes gemacht, also bleibe ich erstmal dabei.
\subsubsection*{Was ist für sie die Zahl 42}
Das ist das Geburtsjahr von meinem Ehemann.2

}{Heiko Carrasco \& Jannis Blüml}

\vfill
\bildmitunterschrift{grafik/comics/fun_with_statistics2.png}{width=11cm}{}{http://abstrusegoose.com/55}

    \clearpage

    % Impressum
    \vfill
\addsec{Impressum}
\small

\textbf{Inforz September \the\year} – Zeitschrift der Studierendenvertretung des Fachbereiches Informatik der Technischen Universität Darmstadt.

\vspace{3mm}
Die Redaktion tagt derzeit unregelmäßig. Die Termine werden über die offene Mailingliste inforz-helfer@D120.de bekannt gegeben. Das Inforz ist  im Web unter \url{www.D120.de/inforz/} verfügbar. Interessierte Mitarbeiter*innen sind immer willkommen; siehe  \url{www.D120.de/inforz/mitmachen/}.

\vspace{3mm}
Namentlich gekennzeichnete und anonyme Beiträge geben nicht unbedingt die Meinung der Redaktion wieder. Alle Rechte, insbesondere das der Verfilmung, vorbehalten. \\

%\vfill

\textbf{Redaktionsanschrift:} Inforz, Fachschaft Informatik, Hochschulstraße 10, 64289 Darmstadt\\
\textbf{Webseite:} \url{www.D120.de/inforz/}\\
\textbf{E-Mail:} inforz@D120.de

\vspace{3mm}
\textbf{Redaktionsschluss dieser Ausgabe:} 20. Juni 2017\\
\textbf{Drucklegung dieser Ausgabe:} 29. August 2017 \\
\textbf{V.i.S.d.P.:} Jannis Blüml, Fachschaft Informatik, Hochschulstraße 10, 64289 Darmstadt

\vspace{3mm}
\textbf{Redaktion:}Tobias Huber, Fabian Franke, Jannis Blüml, Heiko Carrasco

\vspace{3mm}
\textbf{Satz:} Fabian Franke mit \LaTeX\\

\vspace{3mm}
\textbf{Vielen Dank an} Claas Völcker und Tobias Otterbein für einzelne Artikel und alle anderen, die zur Fertigstellung dieses Heftes beigetragen haben.

\vspace{3mm}
\textbf{Titelbild:} Fabian Franke \\
\textbf{Rückumschlag:} Titelseite des Inforz Dezember 1977 \\
\textbf{Comics:} \url{xkcd.com}, \url{abstrusegoose.com}, Creative Commons by-nc

\vspace{3mm}
\textbf{Druck:} die druckerei\\
\textbf{Auflage:} 400 Exemplare\\
\textbf{ISSN:} 1614–4295

\vspace{3mm}
\begin{minipage}{0.2\textwidth}
\includegraphics[scale=.3]{grafik/CreativeCommons/cc}~\includegraphics[scale=.3]{grafik/CreativeCommons/by}~\includegraphics[scale=.3]{grafik/CreativeCommons/sa}
\end{minipage}
\begin{minipage}{0.8\textwidth}
Alle nicht anders gekennzeichneten Artikel sind unter einer \texttt{Creative Commons
Attribution-ShareAlike 4.0 International License}
(\url{creativecommons.org/licenses/by-sa/4.0/}) lizensiert..
\end{minipage}

\vspace{3mm}
\includegraphics[height=20mm]{grafik/eule_final_orange} \hspace{8mm}
\includegraphics[height=18mm]{grafik/logo_ohne_rand}

\pagebreak


    % Rückseite
    
\ThisCenterWallPaper{1}{grafik/inforz_rueckseite}

\thispagestyle{empty}
~

\end{document}
